% ============================================================
% §1  Ontology of TGP — matter as the source of space
% ============================================================
\section{Ontology}\label{sec:ontologia}

\subsection{Central postulate}

The Theory of Generated Space (TGP) rests on a single fundamental
ontological claim:

\begin{axiom}[Space from matter]\label{ax:przestrzen}
Space does not exist independently of matter.
\textbf{The presence of space follows from the presence of matter.}
Every form of mass-energy generates a local contribution to a shared
field~$\Phi$, which measures the ``density of generated space'' at a
given location.  Without matter there is no space; without space there
is no matter.  The two are mutually constitutive.
\end{axiom}

This reverses the standard picture in which spacetime is a pre-existing
arena and matter merely ``moves'' on it.  It also distinguishes TGP from
programmes in which space emerges from information or entanglement
(AdS/CFT, ER\,=\,EPR): in TGP space emerges from \emph{matter}.

\begin{remark}[What TGP is \emph{not}]
TGP is \textbf{not} a scalar-tensor theory with an additional field
living on a pre-existing manifold.  The field~$\Phi$ does not ``inhabit''
a prior spacetime --- it \emph{constitutes} that spacetime.
A scalar-tensor formalism may arise only as an approximation in the
late-time, weak-field limit of TGP, not as its foundation.
\end{remark}

\begin{remark}[Three-level ontology of TGP]\label{rem:ontologia}
TGP adopts the following hierarchy:
\begin{description}[style=nextline,nosep]
  \item[Level 0: Relational substrate $\Gamma = (V, E)$]
    Fundamental and non-emergent.  Possesses adjacency relations
    and connectivity, but no metric.  Exists in both sectors
    $\mathcal{S}_0$ and~$\mathcal{S}_1$.
  \item[Level 1: Spatial field $\Phi$ (metric-phase parameter)]
    $\Phi$ is \textbf{not} space or geometry --- it is a
    \emph{metric-phase parameter} of the substrate, measuring the
    degree to which the substrate supports metric relations
    (distances, propagation).
    $\Phi = 0$: the substrate exists but carries no metric
    (non-metric phase).
    $\Phi > 0$: metric phase --- matter endows the substrate
    with the capacity to support geometry.
  \item[Level 2: Effective geometry $g_{\mu\nu}$]
    Emergent: it arises from coarse-graining Level~1.
    Metric, curvature, and causality are properties of the
    configuration of~$\Phi$, not separate entities.
\end{description}
The substrate is fundamental and non-emergent.  Geometry is a state
of that substrate, controlled by~$\Phi$.  The transition
$\Phi = 0 \to \Phi > 0$ does not create space from nothing ---
it ``switches on'' the metric response of an already existing
substrate.  Analogy: freezing water does not create molecules; it
changes their phase.
\end{remark}

% ----------------------------------------------------------
\subsection{One substrate --- many consequences}\label{rem:jeden-substrat}

The three-level hierarchy above derives from a \textbf{single}
fundamental object: the relational substrate $\Gamma = (V,E)$ with
Hamiltonian~$H_\Gamma$ possessing a chiral $\mathbb{Z}_2$ symmetry
(\S\ref{ssec:chiralna}).  Different coarse-graining projections of
this one object yield:

\begin{enumerate}[nosep]
  \item \textbf{Field $\Phi(x)$} --- the metric-phase parameter,
    obtained from amplitude coarse-graining:
    $\Phi \propto \langle\hat{s}^2\rangle_{\mathcal{B}(x)}$.
    It underpins the \emph{entire} gravitational and cosmological
    sector (Sections~\ref{sec:pole}--\ref{sec:formalizm}).
  \item \textbf{Field $\sigma_{ab}(x)$} --- the spatial-correlation
    tensor:
    $\sigma_{ab} \propto
    \langle\hat{s}_i\,\hat{s}_{i+\hat{e}_a}\rangle_{\mathcal{B}}^{\rm TF}$.
    It carries spin-2 and describes tensor gravitational-wave modes
    (\S\ref{ssec:tensor-substrate}).
    It follows from \textbf{the same} $H_\Gamma$ as $\Phi$ and is
    not an additional postulate.
\end{enumerate}

\noindent\textbf{What this is \emph{not}.}\;
The fields $\Phi$ and $\sigma_{ab}$ do not live ``on'' space; they
\emph{constitute} space (Level~0 $\to$ 1 $\to$ 2).  This is not a
scalar-tensor theory: in a scalar-tensor theory two fields are
\emph{postulated} on a pre-existing manifold.  Here both emerge from
one substrate as different projections.

% ----------------------------------------------------------
\subsection{Nothingness $\Nzero$ and the state function $\Ffield$}%
\label{ssec:N0}

\begin{axiom}[Nothingness]\label{ax:N0}
The fundamental reference state is $\Nzero$ --- absolute nothingness.
Formally:
\begin{equation}
    \Nzero := \bigl\{\emptyset,\;\Ffield\bigr\},
    \qquad
    \Ffield\big|_{\Nzero} = 0,
    \qquad
    \Phi\big|_{\Nzero} = 0.
\end{equation}
In the state $\Nzero$ there is no space ($\Phi = 0$), no matter,
no time, no energy.  The only logical property is that the state
function $\Ffield$ equals zero.
\end{axiom}

$\Nzero$ is not empty space --- it is the absence of space.  It is
not the quantum vacuum --- it is the absence of any structure.

\begin{remark}[Substrate in $\Nzero$: terminological clarification]%
\label{rem:N0-substrat-precyzacja}
In TGP the expression \emph{absolute nothingness} means the
\textbf{absence of all physical metric structure}: no space, matter,
time, or propagation.  It does \emph{not} mean absence of the
substrate~$\Gamma$.  The relational substrate is a Level-0 entity:
it exists in both sectors $\mathcal{S}_0$ and~$\mathcal{S}_1$.
The state $\Nzero$ is the non-metric phase of the substrate ---
a state in which the $\Zp$ symmetry is unbroken
($\langle\hat{s}_i\rangle = 0$ for every~$i$), implying
$\Phi = 0$ and absence of a metric.

\begin{center}
\small
\begin{tabular}{@{}p{4.5cm}p{2.6cm}p{4.5cm}@{}}
\toprule
\textbf{Statement} & \textbf{In TGP} & \textbf{Justification} \\
\midrule
``No metric in $\Nzero$'' & True & $\Phi = 0$ \\
``No matter in $\Nzero$'' & True & $\langle\hat{s}_i\rangle = 0$ \\
``No time in $\Nzero$'' & True & no propagation, $c \to \infty$ \\
``No substrate in $\Nzero$'' & \textbf{False} & $\Gamma$ is fundamental \\
\bottomrule
\end{tabular}
\end{center}
\end{remark}

\begin{definition}[State function $\Ffield$]\label{def:F}
The state function $\Ffield$ measures the \textbf{departure from
nothingness}:
\begin{itemize}[nosep]
  \item $\Ffield = 0$: full symmetry, no structure, state $\Nzero$.
  \item $\Ffield > 0$: local symmetry breaking --- matter exists
        together with the space it generates.
\end{itemize}
At the level of the formalism, $\Ffield$ is realised as the
\textbf{absolute state function}
$\Ffield_{\!\mathrm{abs}}(x) = \Phi(x)^2/\PhiZero^2$,
which vanishes exactly in $\Nzero$ ($\Phi = 0$).
\end{definition}

\begin{axiom}[Zero-sum principle]\label{ax:zero}
Viewed from outside, a closed system (the Universe) is
indistinguishable from~$\Nzero$:
\begin{equation}\label{eq:zero-sum}
    \Ffield_{\mathrm{ext}}\bigl(\text{Universe}\bigr) = 0.
\end{equation}
Internally $\Ffield_{\mathrm{int}}(x) > 0$ locally.  The total net
``amount of existence'' is zero --- the Universe is a self-cancelling
fluctuation of nothingness.
\end{axiom}

% ----------------------------------------------------------
\subsection{Two sectors}\label{ssec:sektory}

\begin{definition}[TGP sectors]\label{def:sektory}
The configuration space $\Omega$ splits into two sectors:
\begin{description}[style=nextline,nosep]
  \item[{$\mathcal{S}_0 = \{\sigma : \Ffield[\sigma] = 0\}$
    --- nothingness sector}]
    No matter, no space, no time.
    The state $\Nzero$ belongs to this sector.
  \item[{$\mathcal{S}_1 = \{\sigma : \Ffield[\sigma] > 0\}$
    --- existence sector}]
    Matter is present, space is generated, time exists as a
    relation between different states.
\end{description}
\end{definition}

% ----------------------------------------------------------
\subsection{Instability of $\Nzero$}\label{ssec:niestabilnosc}

\begin{axiom}[Instability of nothingness]\label{ax:P3}
The state $\Ffield = 0$ ($\Nzero$) is unstable.  Its measure in
configuration space vanishes:
\begin{equation}
    \mu\bigl(\{\sigma \in \Omega : \Ffield[\sigma] = 0\}\bigr) = 0.
\end{equation}
Any departure from exact symmetry leads to $\Ffield > 0$ locally,
entailing the simultaneous appearance of matter and the space it
generates.
\end{axiom}

\begin{proposition}[Instability from substrate thermodynamics]%
\label{prop:P3-deriv}
Let the relational substrate $\Gamma = (V, E)$ have a Hamiltonian
$H_\Gamma$ with chiral $\Zp$ symmetry, and let $\mu_\beta$ be the
Gibbs measure
$\mu_\beta(\sigma) \propto e^{-\beta H_\Gamma[\sigma]}$.
In the thermodynamic limit ($|V| \to \infty$) there exists
$\beta_c > 0$ such that
\begin{equation}
    \beta > \beta_c
    \quad\Longrightarrow\quad
    \mu_\beta\bigl(\{\sigma: \Ffield[\sigma] = 0\}\bigr) = 0,
    \qquad
    \mu_\beta\bigl(\{\sigma: \Ffield[\sigma] > 0\}\bigr) = 1.
\end{equation}
The physical substrate (with no external thermostat) corresponds to
$\beta \to \infty$ --- it is always in the spontaneously broken
regime.  Axiom~\ref{ax:P3} is therefore a \textbf{consequence} of
substrate thermodynamics, not an independent postulate.
\end{proposition}

% ----------------------------------------------------------
\subsection{Chiral symmetry of the substrate}%
\label{ssec:chiralna}\label{ax:symmetry}

On the relational substrate $\Gamma = (V, E)$, every node
$i \in V$ carries an \textbf{amplitude operator} $\hat{s}_i$ ---
a scalar describing the local degree of excitation.

\begin{proposition}[$\Zp$ symmetry from microdynamics]\label{prop:Z2}
The substrate Hamiltonian
\begin{equation}\label{eq:H-Gamma}
    H_\Gamma = \sum_i \biggl[
      \frac{\hat{\pi}_i^2}{2\mu}
    + \frac{m_0^2}{2}\hat{s}_i^2
    + \frac{\lambda_0}{4}\hat{s}_i^4
    \biggr]
    - J\!\sum_{\langle ij\rangle} A_{ij}\,\hat{s}_i\hat{s}_j
\end{equation}
is invariant under \textbf{chiral inversion}:
\begin{equation}\label{eq:Z2-chiral}
    \hat{s}_i \;\mapsto\; -\hat{s}_i
    \qquad \forall\, i \in V.
\end{equation}
This discrete $\Zp$ symmetry is the fundamental symmetry of the
substrate.  The state function $\Ffield$ measures the degree of its
breaking.
\end{proposition}

\begin{corollary}[Relation $\Ffield \leftrightarrow
  \langle\hat{s}\rangle$]\label{cor:F-s}
In the mean-field approximation:
\begin{equation}\label{eq:F-s-relation}
    \Ffield_i[\Omega]
    = \langle\Omega|\hat{s}_i|\Omega\rangle^2.
\end{equation}
$\Ffield = 0 \;\Leftrightarrow\;
\langle\hat{s}_i\rangle = 0$
(unbroken $\Zp$: no structure, nothingness).
$\Ffield > 0 \;\Leftrightarrow\;$ the substrate has ``chosen a side''
--- spontaneous breaking of~$\Zp$ simultaneously generates matter
and space.
\end{corollary}
